\section{Related Work}
\paragraph{Synthetic data with functional environments.} 
Because gathering human-annotated interaction data is notoriously resource-intensive, recent research trends have shifted toward LLM-driven generation. A variety of contemporary works now utilize LLMs to automatically synthesize agentic trajectories at scale~\citep{toolllm, agenttuning, apigen, xu2026controllable, lu2026firefly}. To ensure data fidelity, these approaches execute generated tool calls against live API endpoints or fully implemented sandboxes. While effective for generating grounded trajectories, these methods are fundamentally constrained by their dependence on functional environments. Adapting these pipelines to new APIs requires building bespoke backend infrastructures and painstakingly sourcing the diverse state data required to populate them. Consequently, this overhead severely limits the speed and breadth with which these methods can scale to new API domains.

\vspace{-0.05in}
\paragraph{LLM as an API simulator.}
To bypass the bottleneck of executable backends, several existing works employ LLMs to simulate API responses. Early works such as  ToolAlpaca~\citep{toolalpaca} and ToolACE~\citep{toolace}, and recent efforts like Nemotron~\citep{chandiramani2026nemotron}, synthesize multi-turn data by prompting LLMs to generate mock execution results. However, these methods rely on the LLM's implicit conversational coherence, which leaves them vulnerable to state drift and hallucinated identifiers in long, complex trajectories. Other approaches tailor simulation to specific use cases: ToolEmu~\citep{ruan2024identifying} emulates outputs for safety evaluation, while StableToolBench~\citep{stabletoolbench} uses a GPT-4 simulator as a fallback for the offline evaluation of read APIs. The closest work to ours is Simia~\citep{li2025simulating}, which utilizes simulated environments for agent training by augmenting seed trajectories for SFT and replacing real environments with simulated feedback during RL.

\textsc{ESAT} differs from prior simulation pipelines in four key ways. First, while most simulators focus on read APIs~\citep{toolalpaca, stabletoolbench}, \textsc{ESAT} supports state-changing write APIs, enabling interactions within dynamic, stateful environments. Second, unlike methods requiring pre-compiled API pools~\citep{toolace}, existing benchmarks~\citep{chandiramani2026nemotron}, or seed trajectories~\citep{li2025simulating}, \textsc{ESAT} generates both tasks and trajectories from scratch using only API specifications, making it seamlessly adaptable to new API domains. Third, instead of generating a full trajectory in a single LLM pass as in~\citet{li2025simulating}, \textsc{ESAT} uses explicit, step-by-step interactions between separate agent and simulator LLMs to generate multi-step trajectories. Fourth, to ensure state consistency across these long trajectories, the \textsc{ESAT} simulator validates each simulated response against the full interaction history, thereby acting as a persistent digital world model.